\unnsection{Историческая справка}

    Уравнение конвекции--реакции--диффузии --- это параболическое дифференциальное
    уравнение в частных производных, которое описывает эволюцию концентрации некоторого
    вещества \textit{(примеси, химического реагента, биологической популяции)} под действием
    трёх одновременно протекающих процессов: молекулярной диффузии, конвективного переноса
    с потоком среды и локального источника или стока, обусловленного химической реакцией.

    Историческая особенность данного уравнения состоит в том, что каждый из трёх его
    механизмов имеет собственную, относительно независимую историю и сформировался в рамках
    своей научной традиции --- теории теплопроводности и молекулярной диффузии \textit{(XIX век)},
    химической кинетике \textit{(вторая половина XIX -- начало XX века)} и гидродинамики сплошных
    сред \textit{(с XVIII века)}.

    Объединение этих трёх линий в одно уравнение произошло лишь в первой половине XX века:
    сначала в работе Колмогорова, Петровского и Пискунова 1937 года, где впервые строго
    рассматривалось уравнение, объединяющие диффузию с нелинейной реакцией, а затем в работах
    советской школы теории горения \textit{(Зельдович Я. Б., Франк-Каменецкий Д. А.)}, где
    появилась полная триада <<конвекция--диффузия--реакция>>.

    Дальнейшее изложение организовано в соответствии с этой логикой: сначала рассматривается
    каждая из трёх линий по отдельности, затем --- их объединение, и в заключение --- современные постановки задач.


% =====================================================================
\unnsubsection{Зарождение теории диффузии}
% ---------------------------------------------------------------------
% Объём: ~1 страница.
%
% Ключевые имена и даты:
%   - Жозеф Фурье (1822) — теплопроводность как математическая основа.
%   - Томас Грэм (1828–1833) — первые количественные эксперименты
%     с диффузией газов.
%   - Адольф Фик (1855) — первый и второй законы диффузии по аналогии
%     с законом Фурье; рождение уравнения диффузии.
%   - Роберт Браун (1827) — наблюдение броуновского движения.
%   - Альберт Эйнштейн (1905) — статистическая теория.
%   - Мариан Смолуховский (1906) — независимый вывод.
%   - Жан Перрен (1909, Нобелевская премия 1926) —
%     экспериментальная проверка.
%
% Формулы для включения:
%   - Первый закон Фика: j = -D · ∂c/∂x  (аналогия с Фурье и Омом).
%   - Второй закон Фика (вывод из закона сохранения вещества).
%   - Соотношение Эйнштейна-Смолуховского: <R^2> = 6 D t.
%
% Источники: Mehrer & Stolwijk (разделы 2.1, 2.2, 2.5);
%            оригинальная статья Фика (1855);
%            Wikipedia «Diffusion equation».
% =====================================================================

% TODO: подраздел


% =====================================================================
\unnsubsection{Химическая кинетика и реакционное слагаемое}
% ---------------------------------------------------------------------
% Объём: ~0,5–1 страница.
%
% Ключевые имена и даты:
%   - К. Гульдберг и П. Вааге (1864) — закон действующих масс.
%   - Сванте Аррениус (1889) — температурная зависимость скорости
%     реакции.
%   - Н.Н. Семёнов (Нобелевская премия по химии 1956) — теория
%     цепных реакций; основание советской школы химической кинетики
%     и теории горения.
%   - Д.А. Франк-Каменецкий — мост от химической кинетики к теории
%     горения и к уравнению конвекции-реакции-диффузии.
%
% Формулы для включения:
%   - Закон Аррениуса: k = k_0 · exp(-E / RT).
%   - Краткое пояснение, почему именно эта зависимость делает
%     уравнение конвекции-реакции-диффузии существенно нелинейным.
%
% Источники: Wikipedia «Закон действующих масс»;
%            Mehrer & Stolwijk (раздел 2.4);
%            предисловие к Зельдовичу и др.
% =====================================================================

% TODO: подраздел


% =====================================================================
\unnsubsection{Уравнение Колмогорова--Петровского--Пискунова}
% ---------------------------------------------------------------------
% Объём: ~1–1,5 страницы. Это исторический центр раздела.
%
% Ключевые имена и даты:
%   - А.Н. Колмогоров, И.Г. Петровский, Н.С. Пискунов (1937) —
%     уравнение КПП в контексте биологической задачи о распространении
%     доминантного гена в популяции.
%   - Р.А. Фишер (1937) — параллельная независимая работа о волне
%     преимущественного гена.
%   - Современное название: уравнение Фишера-КПП.
%
% Формулы для включения (из оригинальной статьи КПП-1937):
%   - Исходное уравнение:
%       ∂v/∂t = k · ∂^2 v/∂x^2 + F(v),
%       при F(0) = F(1) = 0,   F(v) > 0 на (0, 1),
%       F'(0) = α > 0,         F'(v) < α на (0, 1).
%   - Скорость предельной бегущей волны:  λ_0 = 2 √(k α).
%   - Уравнение для предельной формы кривой плотности:
%       λ_0 · dv/dx = k · d^2 v/dx^2 + F(v).
%
% Источники: оригинал КПП-1937 (КПП-1937.pdf в папке источников);
%            Wikipedia «Fisher's equation».
% =====================================================================

% TODO: подраздел


% =====================================================================
\unnsubsection{Рождение полного уравнения}
% ---------------------------------------------------------------------
% Объём: ~1 страница. Кульминация раздела — момент, когда «триада»
% (конвекция + диффузия + реакция) собирается воедино в осмысленной
% физической задаче.
%
% Ключевые имена и даты:
%   - Я.Б. Зельдович и Д.А. Франк-Каменецкий (1930–40-е) — теория
%     ламинарного пламени; уравнения распространения пламени, в которых
%     впервые присутствуют все три слагаемых.
%   - Л.Д. Ландау и G. Darrieus — гидродинамическая неустойчивость
%     плоского фронта пламени.
%   - Гидродинамическая основа конвективного слагаемого:
%     Эйлер, Навье-Стокс (упомянуть кратко).
%
% Формулы для включения (из §4 главы 1 монографии Зельдовича и др.):
%   ρ u c · dT/dx = d/dx(λ · dT/dx) + Q · W(a, T)
%   ρ u   · da/dx = d/dx(ρ D · da/dx) − W(a, T)
%
% Здесь явно видны:
%   - конвективные слагаемые ρ u c · d/dx и ρ u · d/dx;
%   - диффузионные слагаемые d/dx(λ · d/dx) и d/dx(ρ D · d/dx);
%   - реакционные слагаемые Q · W и W со скоростью реакции W(a, T),
%     содержащей закон Аррениуса.
%
% При желании — формула нормальной скорости пламени u_n.
%
% Источники: Зельдович и др. (предисловие; глава 1, §4).
% =====================================================================

% TODO: подраздел


% =====================================================================
\unnsubsection{Современные постановки и области применения}
% ---------------------------------------------------------------------
% Объём: ~0,5–1 страница. Плавный переход к разделу
% «Постановка задачи».
%
% Направления применения:
%   - Перенос примеси в атмосфере и водоёмах (школа Г.И. Марчука) —
%     прямая связь с темой курсовой.
%   - Полупроводниковая физика: drift-diffusion equation для
%     концентраций электронов и дырок.
%   - Реакционно-диффузионные системы в биологии: А. Тьюринг (1952) —
%     морфогенез, образование пятен и полос.
%   - Гидрогеология, химические реакторы, теплоперенос в твёрдых телах.
%
% Источники: Wikipedia «Convection-diffusion equation»;
%            Г.И. Марчук «Математическое моделирование в проблеме
%            окружающей среды».
% =====================================================================

% TODO: подраздел
